\chapter[Band 26: Bäume]{Band 26 auf einen Blick}
\noindent\textbf{Bäume}\par\medskip
Vorausgesetzt sind die vier Graphbände und endliche Folgen.
\(\mathcal P_{V,E}(u,v)\) besteht aus Paaren aus Kantenlänge und
injektiver Knotenfolge mit Endpunkten \(u,v\). Bäume müssen nicht endlich sein.

\begin{BandUebersicht}
\textbf{Endliche einfache Verbindungswege}
\UebFormel{\begin{aligned}
&\mathcal P_{V,E}(u,v)=\\
&\{(n,p)\mid n\in\mathbb N,\\
&\quad p\in\powerset(\mathbb N\times V),\\
&\quad\DirectedPath V E p n,\\
&\quad p(0)=u,\ p(n)=v\}.
\end{aligned}}
\UebQuellen{\FormulaRefAuto{FiniteSimplePathFamilyDef}[def]}
&
\textbf{Folgen und Endpunkte sind typisiert}
\UebFormel{\begin{aligned}
&(n,p)\in\mathcal P_{V,E}(u,v),\
 \GraphStruct{V,E}\ \vdash\\
&\quad p:\NatSeg{n+1}\to V,\qquad u,v\in V.
\end{aligned}}
\UebQuellen{\FormulaRefAuto{FiniteSimplePathSequence}[thm];
\FormulaRefAuto{FiniteSimplePathEndpointTyping}[thm]} \\
\addlinespace[.8ex]

\textbf{Wald und Baum}
Ein Wald ist schlicht und besitzt höchstens einen einfachen
Weg zwischen je zwei Knoten. Ein Baum ist schlicht, nichtleer und erfüllt\UebFormel{\begin{aligned}
&\forall u,v\in V\
 \exists!z\ (z\in\mathcal P_{V,E}(u,v)).
\end{aligned}}
\UebQuellen{\FormulaRefAuto{ForestDef}[def];
\FormulaRefAuto{TreeDef}[def]}
&
\textbf{Zusammenhang und Wegeindeutigkeit}
\UebFormel{\begin{aligned}
&\TreeStruct{V,E}\ \vdash\
 \ConnectedGraphStruct{V,E},\\
&\TreeStruct{V,E},\ u,v\in V,\\
&a,b\in\mathcal P_{V,E}(u,v)\ \vdash\ a=b.
\end{aligned}}
\UebQuellen{\FormulaRefAuto{TreeConnected}[thm];
\FormulaRefAuto{TreePathUnique}[thm]} \\
\addlinespace[.8ex]

\textbf{Verwurzelter Baum}
\UebFormel{\begin{aligned}
&\RootedTreeStruct{V,E,r}\ \leftrightarrow\\
&\qquad\TreeStruct{V,E}\land r\in V.
\end{aligned}}
\UebQuellen{\FormulaRefAuto{RootedTreeDef}[def]}
&
\textbf{Eindeutige Wurzelwege genügen}
\UebFormel{\begin{aligned}
&\RootedTreeStruct{V,E,r}\ \eqvdash\\
&\quad\SimpleUndirectedGraphStruct{V,E}\land r\in V\\
&\quad{}\land\forall x\in V\
 \exists!z\ (z\in\mathcal P_{V,E}(r,x)).
\end{aligned}}
\UebQuellen{\FormulaRefAuto{RootedTreeIffUniqueRootPaths}[thm]} \\
\addlinespace[.8ex]

\textbf{Eltern und Kinder}
Im verwurzelten Baum ist \(a\) Elternknoten von \(x\ne r\),
wenn es auf dessen Wurzelweg unmittelbar vor \(x\) steht:\UebFormel{\begin{aligned}
&\exists k\in\mathbb N\,\exists p\\
&\quad((k+1,p)\in\mathcal P_{V,E}(r,x)\\
&\qquad\land p(k)=a).
\end{aligned}}Kinder sind durch die umgekehrte Elternrelation definiert.
\UebQuellen{\FormulaRefAuto{TreeParentDef}[def];
\FormulaRefAuto{TreeChildDef}[def];
\FormulaRefAuto{TreeChildrenDef}[def]}
&
\textbf{Genau ein Elternknoten außerhalb der Wurzel}
\UebFormel{\begin{aligned}
&\RootedTreeStruct{V,E,r},\
 x\in V\setminus\{r\}\ \vdash\\
&\qquad\exists!a\in V\
 \TreeParentRel V E r a x.
\end{aligned}}Die Wurzel hat keinen Elternknoten.
\UebQuellen{\FormulaRefAuto{TreeParentUnique}[thm];
\FormulaRefAuto{RootHasNoParent}[thm]} \\
\addlinespace[.8ex]

\textbf{Blätter, innere Knoten und Vorfahren}
Im verwurzelten Baum sei \(K(x)=\TreeChildren V E r x\).\UebFormel{\begin{aligned}
&\operatorname{Blatt}(x)\leftrightarrow
      (x\in V\land K(x)=\varnothing),\\
&\operatorname{Innen}(x)\leftrightarrow
      (x\in V\land K(x)\ne\varnothing).
\end{aligned}}Ein Vorfahr liegt auf dem Wurzelweg zum betrachteten Knoten.
\UebQuellen{\FormulaRefAuto{TreeLeafInnerDef}[def];
\FormulaRefAuto{TreeAncestorDef}[def]}
&
\textbf{Zerlegung und ausgezeichnete Vorfahren}
Für \(\RootedTreeStruct{V,E,r}\) und \(x\in V\):\UebFormel{\begin{aligned}
&\operatorname{Blatt}(x)\lor\operatorname{Innen}(x),\\
&\operatorname{Blatt}(x)\ \vdash\neg\operatorname{Innen}(x),\\
&\TreeAncestorRel V E r x x,\\
&\TreeAncestorRel V E r r x.
\end{aligned}}
\UebQuellen{\FormulaRefAuto{TreeLeafInnerPartition}[thm];
\FormulaRefAuto{TreeAncestorReflexive}[thm];
\FormulaRefAuto{RootAncestorEveryVertex}[thm]} \\
\addlinespace[.8ex]

\textbf{Binäre und geordnete volle Binärbäume}
Binär bedeutet: höchstens zwei Kinder je Knoten. Voll bedeutet:
null oder genau zwei Kinder. Eine Ordnung benennt auf der Menge \(I\)
der inneren Knoten linkes und rechtes Kind:\UebFormel{\begin{aligned}
&\lambda,\rho:I\to V,\quad
 \lambda(x)\ne\rho(x),\\
&K(x)=\{\lambda(x),\rho(x)\}\quad(x\in I).
\end{aligned}}
\UebQuellen{\FormulaRefAuto{BinaryTreeDef}[def];
\FormulaRefAuto{FullBinaryTreeDef}[def];
\FormulaRefAuto{OrderedFullBinaryTreeDef}[def]}
&
\textbf{Innere Knoten haben zwei verschiedene Kinder}
\UebFormel{\begin{aligned}
&\FullBinaryTreeStruct{V,E,r},\
 \TreeInnerPred V E r x\ \vdash\\
&\quad\exists a,b\in K(x)\
 (a\ne b\land K(x)=\{a,b\}).
\end{aligned}}Vergessen von \(\lambda,\rho\) liefert einen vollen Binärbaum,
weiter einen binären und einen verwurzelten Baum.
\UebQuellen{\FormulaRefAuto{FullBinaryInnerHasTwoChildren}[thm];
\FormulaRefAuto{OrderedFullBinaryTreeForgetfulChain}[thm]} \\
\addlinespace[.8ex]
\end{BandUebersicht}
